\section{Use explicit interfaces at each boundary}
Keep file bytes, parsed metadata and display words in separate objects. The byte
interfaces accept unaligned source and serialized-output pointers. Framebuffer words
require four-byte alignment; the metadata object requires eight-byte alignment.
Outputs are disjoint from their inputs. Counts, capacities, pointers, saved registers
and return addresses retain all 64 bits.

The parser \texttt{ppm\_parse} receives a byte pointer in \texttt{a0}, its unsigned
count in \texttt{a1} and a metadata pointer in \texttt{a2}. On success it writes sixteen
bytes: a 32-bit width at offset zero, a 32-bit height at offset four, and a 64-bit raster
offset at offset eight. All fields use little-endian memory order. Its return in
\texttt{v1} is the number of pixels. No metadata byte changes on failure. The parser
validates the complete extent, but it need not inspect each sample's numerical value:
every byte value is a valid sample when Maxval is 255.

The decoder \texttt{ppm\_decode} receives the bytes and count in \texttt{a0,a1}, a
framebuffer pointer in \texttt{a2} and its capacity in pixels in \texttt{a3}. It returns
the dimensions in \texttt{v1}: width in bits 31--0, height in bits 63--32. For the
worked image this is \texttt{0x0000000200000002}. The encoder \texttt{ppm\_encode}
accepts a framebuffer pointer, that packed dimension value, a serialized-output pointer
and its capacity in bytes. Its successful \texttt{v1} is the exact file byte count.
The encoder ignores each input word's high byte and emits only its RGB fields. It
writes no trailing terminator or guard byte.

\Needspace{7\baselineskip}
For memory interfaces, \texttt{v0=0} means success, one means malformed input or invalid
dimensions, and two means a bound or capacity failure. Failures return zero in
\texttt{v1}. The parser and decoder reject a byte count above 12544 before accessing
pointers; an empty or otherwise malformed bounded file returns one. The decoder
validates the file before checking framebuffer capacity. The encoder validates both
dimensions before checking output capacity, and both checks precede pixel reads or
output writes. Preserve all bits of \texttt{s0} through \texttt{s7}, \texttt{fp} and
\texttt{sp}; produce no console output. Temporary registers are volatile.

\section{Give file buffers and files different failure promises}
The loader \texttt{ppm\_load} receives a path in \texttt{a0}, a byte buffer in
\texttt{a1}, its capacity in \texttt{a2}, and a metadata pointer in \texttt{a3}.
Capacity must be from one through 12544. An invalid capacity returns two before opening
a path or accessing caller objects. The loader reads to EOF, closes the file, then
validates the complete bytes before publishing metadata. Successful \texttt{v1} is the
file byte count. A format error returns one; excess input returns two; an open or read
failure returns three. Its byte buffer may contain a prefix on failure, but metadata
remains unchanged. This permits buffering while retaining the later display boundary.

The saver \texttt{ppm\_save} receives a path, serialized bytes and their count in
\texttt{a0,a1,a2}. It validates the file before opening the destination in truncate
mode. A zero count returns one and an over-limit count returns two before pointer
access. An open or write failure returns three. Failures return zero in
\texttt{v1}; successful \texttt{v1} is the saved count. A write failure can leave a
partial output file. Closing a descriptor does not roll those bytes back.

Both file routines close every successful open on every return path and preserve the
same saved registers as the memory routines. Paths are NUL-terminated byte strings.
Supply the extra readable guard byte described in the syscall discussion after any
serialized file passed to the saver. It is outside the file's declared count.
